\begin{abstract}
The emergence of 1-bit Large Language Models (LLMs), exemplified by research efforts like BitNet~\cite{bitnet}, signals a paradigm shift within the field. We present a ternary LLM, \textbf{\text{BitNet b1.58}}, whose weights take values exclusively from the set \{-1, 0, 1\}. This design allows the model to achieve parity with conventional full-precision (FP16 or BF16) Transformer LLMs of equivalent parameter count and training volume, maintaining competitive results in perplexity and downstream task effectiveness. Notably, the 1.58-bit LLM achieves these outcomes with substantial improvements in latency, memory footprint, computational throughput, and power efficiency. On a broader scale, this development establishes a revised scaling law and a novel methodological foundation for building next-generation LLMs that offer both high performance and reduced operational costs. Additionally, it unlocks fresh possibilities in computing models, inviting exploration into bespoke hardware architectures tailored for 1-bit LLM workloads.
\end{abstract}

\section{The Era of 1-bit LLMs}

Artificial intelligence has recently experienced dramatic advances with the escalating scale and sophistication of Large Language Models (LLMs). These models excel across diverse natural language processing applications, yet the rapid expansion of their parameter counts presents difficulties for real-world deployment and highlights growing environmental and financial concerns owing to escalating energy demands.

A prominent strategy for mitigating such constraints is post-training quantization, which compresses models to lower precision for inference~\cite{smoothquant, gptq, quip, quip_sharp}. By reducing the bit-width of both weights and activations, this approach drastically cuts the memory and computation overhead required by LLMs. The prevailing movement within the community is toward ever lower bit representations, progressing from traditional 16-bit to contemporary 4-bit configurations~\cite{gptq, awq}. Despite its widespread adoption across industry applications, post-training quantization is not optimal and exhibits intrinsic limitations.

Recent advancements in 1-bit model architectures, exemplified by BitNet~\cite{bitnet}, indicate a fruitful approach for reducing the resource demands of LLMs without compromising performance. Conventional LLMs use 16-bit floating point representations (such as FP16 or BF16), and most of their operations revolve around matrix multiplications, incurring considerable computational costs due to floating-point additions and multiplications. In comparison, BitNet executes matrix multiplications relying solely on integer additions, offering substantial energy efficiency improvements over traditional LLMs. Since power constraints often define the ceiling for computing performance on modern chips, this reduction in energy usage directly contributes to computational speedups.

Beyond raw computation, the overhead associated with moving model weights from DRAM into the on-chip accelerator memory (e.g., SRAM) can be substantial during inference. Strategies aiming to increase SRAM capacity may boost throughput, yet come with a far greater expense than expanding DRAM. Relative to their full-precision counterparts, 1-bit LLMs demand significantly less memory both in terms of storage and bandwidth. This reduced requirement allows for more cost- and time-efficient transfer of model weights from DRAM, enhancing inference speed and overall efficiency.

Within this study, we present a notable 1-bit LLM extension named \textbf{\text{BitNet b1.58}}, where each parameter is represented as a ternary value drawn from \{-1, 0, 1\}. This formulation extends the original 1-bit BitNet by incorporating an explicit 0, equating to 1.58 bits in binary terms. \text{BitNet b1.58} preserves all core advantages of the earlier 1-bit BitNet, most notably its computation approach that nearly eliminates multiplication in matrix operations, enabling aggressive optimization. Its energy footprint remains equivalent to the 1-bit BitNet, while its memory efficiency, throughput, and latency significantly surpass FP16 LLMs. Additionally, \text{BitNet b1.58} introduces two noteworthy enhancements. First, it has increased modeling flexibility, attributable to the introduction of 0 in the parameter set, which enables explicit feature filtering—substantially improving the effectiveness of 1-bit LLMs. Second, experimental results demonstrate that \text{BitNet b1.58} achieves parity with full-precision (FP16) baselines on both perplexity and downstream tasks, beginning at the 3B scale, when matched for factors such as architecture and training regimen.

\section{BitNet b1.58}

\text{BitNet b1.58} builds upon the foundational BitNet architecture, characterized as a Transformer variant where the conventional \emph{nn.Linear} layers are substituted with \emph{BitLinear} layers. This model is trained entirely from the beginning, utilizing 1.58-bit precision for its weights and 8-bit for activations. In contrast to the initial BitNet version, this iteration incorporates several adjustments, which we detail below.

\paragraph{Quantization Function.} In order to restrict the weights to values in the set $\{-1, 0, +1\}$, an \emph{absmean} quantization approach is employed. This method begins by normalizing the weight matrix with its mean absolute value, after which each entry is discretized to the closest integer within the allowed set:

\begin{equation}
    \widetilde{W} = \mathrm{RoundClip}(\frac{W}{\gamma + \epsilon}, -1, 1),
\end{equation}

\begin{equation}
    \mathrm{RoundClip}(x, a, b) = \max(a, \min(b, \mathrm{round}(x))),
\end{equation}

\begin{equation}
    \gamma = \frac{1}{nm}\sum_{ij} |W_{ij}|.
\end{equation}

For activations, the quantization follows the same strategy implemented in BitNet; however, we omit the rescaling of activations to the interval $[0, Q_b]$ prior to applying nonlinearities. Instead, activations are rescaled to $[-Q_b, Q_b]$ for each token independently, thus eliminating the need for zero-point quantization. This alternative not only simplifies implementation and system-level deployment but, according to our experiments, has an insignificant impact on model performance.

\paragraph{LLaMA-alike Components.}

LLaMA~\cite{llama, llama2} has become a standard reference model for open-source large language models. Aligning with this community, \text{BitNet b1.58} adopts several LLaMA-inspired features. It integrates RMSNorm~\cite{rmsnorm}, SwiGLU~\cite{swiglu}, rotary position embeddings~\cite{rope}, and omits bias terms throughout. This compatibility facilitates straightforward integration with prevalent open-source libraries such as Huggingface, vLLM~\cite{vllm}, and llama.cpp\footnote{\url{https://github.com/ggerganov/llama.cpp}}, minimizing adaptation overhead.

\section{Results}

We evaluated \text{BitNet b1.58} alongside our implementation of the FP16 \text{LLaMA LLM} across a variety of model scales. Both models were pre-trained using the RedPajama dataset~\cite{redpajama} for a total of 100 billion tokens, ensuring comparability. Their zero-shot abilities were assessed across several language benchmarks, such as ARC-Easy~\cite{arc}, ARC-Challenge~\cite{arc}, Hellaswag~\cite{hellaswag}, Winogrande~\cite{winoGrande}, PIQA~\cite{piqa}, OpenbookQA~\cite{openbookqa}, and BoolQ~\cite{boolq}. Additionally, we included validation perplexity scores for WikiText2~\cite{wikitext2} and C4~\cite{c4}.

In addition to performance metrics, we analyzed the GPU memory consumption and latency during inference for both \text{LLaMA LLM} and \text{BitNet b1.58}. The measurements were obtained using the FasterTransformer\footnote{\url{https://github.com/NVIDIA/FasterTransformer}} framework, optimized for minimizing latency on GPU-based LLM inference. For \text{BitNet b1.58}, we also incorporated Ladder’s~\cite{ladder} 2-bit kernel. Our primary latency metric was the duration required to generate each output token, reflecting the dominant inference cost.

Table~\ref{tab:ppl} presents both the perplexity and the associated computational costs for \text{BitNet b1.58} and \text{LLaMA LLM}. Notably, at the 3B parameter scale, \text{BitNet b1.58} reaches the perplexity levels of its full precision counterpart, while delivering a 2.71$\times$ speedup and consuming 3.55$\times$ less GPU memory. Furthermore, at the 3.9B scale, \text{BitNet b1.58} is 2.4$\times$ faster, reduces memory usage by 3.32$\times$, and outperforms the 3B \text{LLaMA LLM} model.

Table~\ref{tab:zero-shot} details the zero-shot accuracy across various tasks. The evaluation followed the \emph{lm-evaluation-harness}\footnote{\url{https://github.com/EleutherAI/lm-evaluation-harness}} workflow. The results indicate that as the model size increases, the accuracy disparity between \text{BitNet b1.58} and \text{LLaMA LLM} diminishes. Importantly, from the 3B size onward, \text{BitNet b1.58} achieves parity with the full precision baseline. Echoing the perplexity findings, task-based evaluations show that \text{BitNet b1.58} at 3.9B outperforms the 3B \text{LLaMA LLM} while maintaining lower latency and memory consumption. These findings establish \text{BitNet b1.58} as a Pareto improvement over current leading LLM architectures.

\paragraph{Memory and Latency}

We extended our scaling experiments to include model sizes of 7B, 13B, and 70B, and assessed the corresponding computational costs. As depicted in Figure~\ref{fig:latency-memory-energy}, both latency and memory usage display clear scaling trends: speed improvements become more significant with larger model sizes. Notably, the \text{BitNet b1.58} model at 70B parameters achieves a 4.1$\times$ speed advantage over the \text{LLaMA LLM} baseline. This improvement stems from the fact that the computational demands of the \emph{nn.Linear} layer increase proportionally with the scale of the model. Memory requirements show analogous behavior; as embeddings retain full precision, their memory footprint represents a diminishing fraction in larger models. Latency and memory benchmarks were obtained using a 2-bit kernel implementation, suggesting that further efficiency gains remain attainable through kernel optimization.

\paragraph{Energy}

We further estimated the energy requirements for arithmetic operations in both the \text{BitNet b1.58} and \text{LLaMA LLM} models, centering our analysis on matrix multiplications as the principal source of computation in LLMs. The breakdown of energy expenditure, presented in Figure~\ref{fig:energy}, reveals that INT8 addition operations constitute the bulk of \text{BitNet b1.58}'s workload, whereas \text{LLaMA LLM} predominantly relies on FP16 addition and multiplication. Employing the energy consumption model referenced in~\cite{energycost, pokebnn}, we find that \text{BitNet b1.58} reduces arithmetic operation energy for matrix multiplication by a factor of 71.4 on 7nm hardware. Additionally, we report comprehensive end-to-end energy measurements using input sequences of 512 tokens. Our findings indicate that the energy efficiency of \text{BitNet b1.58}, relative to the FP16 \text{LLaMA LLM} baseline, improves with increasing model size. This outcome is attributable to the rising contribution of \emph{nn.Linear} to total cost as models scale, while other overhead components become proportionally less significant.

\paragraph{Throughput}

We evaluate the throughput of \text{BitNet b1.58} alongside the \text{LLaMA LLM} comprising 70B parameters across two 80GB A100 GPUs. Pipeline parallelism~\cite{gpipe} is employed to make it feasible to run \text{LLaMA LLM} 70B on the given hardware. By increasing the batch size incrementally until reaching the GPUs' memory capacity, with a sequence length fixed at 512, we find that, as shown in Table~\ref{tab:throughput}, the \text{BitNet b1.58} 70B supports batch sizes up to 11 times larger than those supported by the \text{LLaMA LLM}, resulting in an 8.9-fold increase in throughput.

\textbf{\text{BitNet b1.58} establishes a novel scaling law relating model efficiency to inference cost and model capacity.} Drawing from the findings illustrated in Figure~\ref{fig:latency-memory-energy} and \ref{fig:energy}, we can map the following correspondences between models at 1.58-bit and their 16-bit FP16 counterparts:

\begin{itemize}
    \item 13B BitNet b1.58 surpasses the 3B FP16 LLM in efficiency for latency, memory consumption, and energy requirements.
    \item 30B BitNet b1.58 demonstrates better efficiency in latency, memory usage, and energy demand compared to the 7B FP16 LLM.
    \item 70B BitNet b1.58 is more efficient across latency, memory footprint, and energy consumption than the 13B FP16 LLM.
\end{itemize}

\paragraph{Training with 2T Tokens}

The total count of tokens utilized during training is a critical parameter for LLM performance. To assess how well \text{BitNet b1.58} scales with respect to training token volume, we trained a \text{BitNet b1.58} variant using 2T tokens, employing the same dataset construction approach as StableLM-3B~\cite{StableLM-3B-4E1T}, the current leading open-source 3B model. Both models were then tested on a benchmark suite that includes Winogrande~\cite{winoGrande}, PIQA~\cite{piqa}, SciQ~\cite{sciq}, LAMBADA~\cite{lambada}, and ARC-easy~\cite{arc}, reporting zero-shot accuracy as shown in Table~\ref{tab:2t}. For those tasks assessed by both accuracy and normalized accuracy, we averaged the two scores. The results for StableLM 3b at 2T tokens were directly taken from its official report. Our experimental results indicate that \text{BitNet b1.58} delivers superior outcomes across all tasks, suggesting robust generalization ability for 1.58-bit LLMs as well.

\section{Discussion and Future Work}

\textbf{1-bit Mixture-of-Experts (MoE) LLMs}

Mixture-of-Experts (MoE) architectures offer a resource-efficient solution for large language models (LLMs), notably lowering computational FLOPs. However, their practical use is hindered by substantial memory demands and inter-chip communication bottlenecks. Employing 1.58-bit LLMs provides remedies for these issues. With a smaller memory footprint, fewer devices are necessary to host MoE models, and the overhead involved in transmitting activations across chips or networked systems can be dramatically diminished. Ideally, fitting entire models on a single chip would completely eliminate inter-device transfer costs.

\textbf{Native Support of Long Sequence in LLMs}

Managing extended sequence lengths has become increasingly vital for LLMs. The primary obstacle lies in the considerable memory consumption of the KV caches during long sequence inference. \text{BitNet b1.58} marks a pivotal advance by halving the activation precision from 16 bits to 8 bits, thereby permitting a doubling of context length under identical hardware constraints. Further improvements are possible, as activations in 1.58-bit LLMs can potentially be compressed losslessly to 4 bits or below, which remains an open avenue for exploration.

\textbf{LLMs on Edge and Mobile}

Adopting 1.58-bit LLMs can lead to substantial enhancements for running language models on edge and mobile platforms, which are typically constrained by limited computational resources and memory. The efficient memory and power usage characteristics of 1.58-bit LLMs make it feasible to deploy them in such environments, paving the way for a broader spectrum of applications that were not previously attainable. This development stands to notably advance the functionality of edge and mobile hardware and facilitate innovative uses of LLMs. Furthermore, 1.58-bit LLMs are particularly suitable for CPU-based systems—the predominant architecture in edge and mobile devices—making \text{BitNet b1.58} an effective and performant choice for such contexts.

\textbf{New Hardware for 1-bit LLMs}

Emerging solutions such as Groq\footnote{\url{https://groq.com/}} illustrate the possibilities for specialized hardware (e.g., LPUs) designed to accelerate LLM workloads. As a natural progression, we advocate for focused research and development of new hardware and system designs tailored to the distinct requirements of 1-bit LLMs, leveraging the innovative computational framework established by BitNet~\cite{bitnet}.